\documentclass[11pt]{article}

\usepackage[margin=2.2cm]{geometry}
\usepackage{booktabs}
\usepackage{array}
\usepackage{tabularx}
\usepackage{longtable}
\usepackage{amsmath}
\usepackage[colorlinks,breaklinks]{hyperref}

\begin{document}

\title{Supplementary Methods for ``DentalClaw: A Modular Workflow Platform for Standardized and Auditable Dental Imaging AI Experiments''}
\author{The Authors}
\date{}
\maketitle

\section{Workflow-level outcomes}

Table~S1 summarizes end-to-end workflow completion, preparation findings, and stage timings for the two representative tasks, together with a manually assembled script workflow executed for comparison. Stage times are elapsed wall-clock time on the same NVIDIA RTX 3090 workstation. The manual arm was performed by the AI researcher on the study team following the same predefined workflow without additional hyperparameter tuning or model selection; operator background and timing rules are reported in the manual comparison protocol below, and the comparison is a workflow demonstration rather than a controlled human-factor experiment.

\begin{table}[!ht]
\centering
\small
\setlength{\tabcolsep}{4pt}
\caption{End-to-end workflow completion, preparation findings, and stage timings. Stage times are elapsed wall-clock time on the same NVIDIA RTX 3090 workstation.}
\label{tab:supp_workflow}
\begin{tabular}{l c c p{4.6cm} p{4.8cm}}
\toprule
Task & Cases input / eligible & Main preparation findings & DentalClaw (wall-clock) & Manual workflow (wall-clock) \\
\midrule
Panoramic segmentation & 1051 / 968 & Missing masks; image--label mismatch & audit+prep 14 min; train 6 h 52 min; infer+eval+report 31 min & audit+prep 3 h 26 min; train 6 h 58 min; infer+eval+organization 2 h 18 min; outputs: ad hoc logs and files \\
\addlinespace
CBCT segmentation & 63 / 43 & Metal artifacts; FOV anomalies; inconsistent volume usability & audit+prep 21 min; train 19 h 36 min; infer+eval+report 1 h 44 min & audit+prep 6 h 12 min; train 19 h 49 min; infer+eval+organization 3 h 11 min; outputs: ad hoc logs and files \\
\bottomrule
\end{tabular}
\end{table}

\section{Training configuration}

\subsection{TDD 32-class panoramic segmentation configuration trajectory}

The 32-class configuration trajectory was executed under the same audited dataset, predefined train/validation/test split, and evaluation pipeline. The baseline was a default U-Mamba model; each subsequent experiment added one or more reusable configuration skills (Table~S2). Mean Dice and mean HD95 are reported at the image level.

\begin{table}[!ht]
\centering
\small
\caption{Configuration trajectory executed by DentalClaw on a 32-class panoramic dental X-ray dataset.}
\label{tab:supp_tdd_trajectory}
\begin{tabular}{l l c c}
\toprule
Experiment & Key modification & Mean Dice $\uparrow$ & Mean HD95 $\downarrow$ (mm) \\
\midrule
Baseline & Default U-Mamba (no extra optimisations) & 0.6163 & 83.15 \\
Exp1 & Gaussian noise + gamma correction (training-time) & 0.6666 & 50.42 \\
Exp2 & Cosine annealing learning-rate schedule & 0.6483 & 56.90 \\
Exp3 & AdamW optimizer & 0.6592 & 58.74 \\
Exp4 & Test-time augmentation (mirror) & 0.6208 & 82.78 \\
Exp5 & Basic combination: AdamW + TTA (mirror) & 0.6684 & 48.52 \\
\textbf{Exp6} & \textbf{Best combination: AdamW + Cosine + TTA} & \textbf{0.6937} & \textbf{38.54} \\
\bottomrule
\end{tabular}
\end{table}

\subsection{User-selected configuration comparison}

As a documented reference rather than a formal human-factor experiment, three operators with different AI backgrounds (Student A: AI novice; Student B: master's-level researcher; AI researcher: practitioner in the field) manually selected configurations for the same 32-class task under the same split and evaluation pipeline (Table~S3).

\begin{table}[!ht]
\centering
\small
\caption{User-selected vs.\ DentalClaw skill-based configurations under the same dental X-ray task.}
\label{tab:supp_user_comparison}
\begin{tabular}{l l c c}
\toprule
Source & Configuration type & Mean Dice $\uparrow$ & Mean HD95 $\downarrow$ (mm) \\
\midrule
Student A & User-selected & 0.6283 & 82.28 \\
Student B & User-selected & 0.6546 & 56.37 \\
AI researcher & User-selected & 0.6883 & 49.21 \\
\textbf{DentalClaw} & \textbf{Skill-based} & \textbf{0.6937} & \textbf{38.54} \\
\bottomrule
\end{tabular}
\end{table}

\subsection{CBCT 35-class segmentation configuration trajectory}

The CBCT task involved 35 anatomical and restorative classes (teeth 1--32 and restorations 33--35). DentalClaw coordinated three predefined segmentation modules (nnU-Net, SwinUNETR, U-Mamba) with multi-model ensemble and post-processing skills under one traceable workflow (Table~S4).

\begin{table}[!ht]
\centering
\small
\caption{Configuration trajectory executed by DentalClaw on a 35-class CBCT dataset.}
\label{tab:supp_cbct_trajectory}
\begin{tabular}{l l c c c c}
\toprule
& & \multicolumn{2}{c}{Teeth (1--32)} & \multicolumn{2}{c}{Restorations (33--35)} \\
\cmidrule(lr){3-4} \cmidrule(lr){5-6}
Experiment & Key modification & Mean Dice $\uparrow$ & HD95 $\downarrow$ (mm) & Mean Dice $\uparrow$ & HD95 $\downarrow$ (mm) \\
\midrule
Baseline & Default nnU-Net (CNN) & 0.7038 & 10.386 & 0.6588 & 16.484 \\
Exp1 & SwinUNETR (Transformer) & 0.6508 & 15.633 & 0.7371 & 9.625 \\
Exp2 & U-MambaBot (SSM) & 0.6648 & 11.886 & 0.7279 & 6.122 \\
Exp3 & Ensemble (majority voting) & 0.7189 & 11.231 & 0.7508 & 5.551 \\
Exp4 & Ensemble (mean-softmax) & 0.7156 & 11.080 & 0.7536 & 5.888 \\
\textbf{Exp5} & \textbf{Best: ensemble + TTA + post-processing} & \textbf{0.8058} & \textbf{4.528} & \textbf{0.7653} & \textbf{5.172} \\
\bottomrule
\end{tabular}
\end{table}

\subsection{TDD binary segmentation route (platform registry)}

The platform registry route \texttt{tdd\_2d\_segmentation\_infer\_report} uses an existing nnU-Net v2 checkpoint in the 2D full-resolution configuration for binary tooth segmentation (one class), with mean Dice as the primary metric. The nnU-Net dataset layout contains 900 training and 100 test cases, generated by deterministic case-level splitting with a fixed random seed (42). Five trainer variants (default, lower learning rate, higher weight decay, and swapped scheduler configurations) were executed through the platform's configuration-search skill; fold-all validation mean Dice ranged from 0.9312 to 0.9331, with the best variant reaching 0.9331.

\subsection{Training configuration details}

All models were trained with nnU-Net v2 using the framework's default training recipe unless stated otherwise: stochastic gradient descent with Nesterov momentum, polynomial learning-rate decay from an initial rate of 0.01, and a combined cross-entropy and soft Dice loss. The 32-class panoramic trajectory used the U-Mamba architecture trained for 100 epochs; the configuration skills modified this base recipe as follows: Gaussian-noise and gamma-correction augmentation at training time (Exp1), a cosine-annealed learning-rate schedule (Exp2), the AdamW optimizer with initial learning rate $1\times10^{-4}$ and weight decay $1\times10^{-4}$ (Exp3), and mirroring test-time augmentation (Exp4). The 2D configuration used training patches of 896$\times$1536 pixels at 1.0\,mm spacing with a batch size of 2 and Z-score intensity normalization; the backbone was a nine-stage plain U-Net with instance normalization, leaky ReLU activations, and 32 base features. The 35-class CBCT models used the 3D full-resolution configuration with patches of 96$\times$160$\times$128 voxels at 0.3\,mm isotropic spacing, a batch size of 2, and Z-score normalization. Ensemble and post-processing skills aggregated single-model predictions as described in Table~S4.

\subsection{Software environment}

All workflow runs used a workstation with an NVIDIA GeForce RTX 3090 GPU. Model execution used nnU-Net v2 under a Python 3.10 environment on Linux. Three-dimensional surface renderings of segmentation masks were generated with 3D Slicer (version 5.10.0).

\section{Private dataset details}

The private panoramic dataset comprised 290 radiographs collected during routine clinical care between 2012 and 2024 under institutional ethical approval (PKUSSIRB-2025,107,012). Images were acquired with a Planmeca dental panoramic X-ray system and stored as 8-bit grayscale JPEG files. The cohort had a mean age of 35.6 years (standard deviation 15.2; range 16--80) and a gender distribution of 51\% male / 49\% female. Clinical experts annotated bounding boxes for four categories: pulp, root canal fillings, fillings, and periapical lesions. Images with poor quality, severe artifacts, incomplete dentition coverage, or uncommon anatomical conditions that could compromise annotation reliability were excluded.

\section{Intent-to-workflow decision evaluation}

\subsection{Intent set and expected outcomes}

The evaluation used a study-defined set of 30 natural-language intents covering four behavior classes (5 standard executable requests, 13 boundary requests without a validated route, 7 ambiguous requests, and 5 trap requests), three dataset groups (TDD, ToothFairy3, and private 2D), and four task families (2D segmentation, 3D segmentation, detection, and classification). Several requests deliberately omitted the task family from the request text. For every intent, the expected platform outcome (supported, executable, and target method) was prespecified by the study team before evaluation on the basis of the offline registry contents and the platform's documented scope. Table~S5 lists the full request texts, categories, expected outcomes, scores, and actual decisions.

\begin{table}[!ht]
\centering
\footnotesize
\caption{Full results of the 30-intent decision evaluation. Ext: non-executable \texttt{agent\_external\_suggestion}; Clar: \texttt{platform\_clarification}; Rej: rejected (\texttt{supported=false}). Route abbreviations: tdd-inf = \texttt{tdd\_2d\_segmentation\_infer\_report}; tf3 = \texttt{toothfairy3\_3d\_segmentation\_infer\_or\_train}; pvt = \texttt{private\_2d\_segmentation\_train}.}
\label{tab:supp_intents}
\begin{tabular}{@{}c p{9.0cm} c c c@{}}
\hline
\# & \textbf{Request (Chinese, as issued)} & \textbf{Category} & \textbf{Score} & \textbf{Decision} \\
\hline
\multicolumn{5}{@{}l@{}}{\textbf{Standard executable requests}} \\
1 & 用已有的 TDD 二值分割模型在测试集上推理，并输出 Dice、IoU、HD95 和叠加图。 & standard & 1.000 & tdd-inf \\
2 & 在 ToothFairy3 CBCT 数据上训练 3D 多类别牙齿分割模型。 & standard & 1.000 & tf3 \\
3 & 在 ToothFairy3 全量 CBCT 上训练 3D 分割模型；如果 QC 发现需人工复核或阻断病例，停止并输出病例清单。 & standard & 1.000 & tf3 \\
4 & 用已有 ToothFairy3 3D 分割模型对验证/测试 CBCT 做推理评估。 & standard & 1.000 & tf3 \\
5 & 用私有 2D 数据 imagesTr/labelsTr 训练牙齿分割模型，并用 imagesVal/labelsVal 验证。 & standard & 1.000 & pvt \\
\hline
\multicolumn{5}{@{}l@{}}{\textbf{Boundary requests (no validated route)}} \\
6 & 用 TDD 全景片训练一个默认牙齿二值分割模型，并保留 10\% 测试集。 & boundary & 0.985 & Ext \\
7 & 把 TDD 的牙齿多边形标注导出为 32 类分割，并训练一个 FDI/32 类牙位分割模型。 & boundary & 0.985 & Ext \\
8 & 把 TDD 的 teeth\_bbox.json 整理成 COCO 牙齿检测数据集，供检测模型训练使用。 & boundary & 0.985 & Ext \\
9 & 基于 TDD bbox 标注训练一个牙齿检测模型，并报告验证集 mAP。 & boundary & 0.985 & Ext \\
10 & 用 TDD 全景片训练一个病例级疾病分类模型。 & boundary & 0.940 & Rej \\
11 & 从 ToothFairy3 3D 标签中导出牙齿检测框数据集。 & boundary & 0.985 & Ext \\
12 & 用 ToothFairy3 训练 3D 牙齿检测模型，并报告检测 mAP。 & boundary & 0.985 & Ext \\
13 & 将 ToothFairy3 病例分成 usable 和 needs\_manual\_review，并输出病例级分类结果。 & boundary & 0.985 & Ext \\
14 & 对私有 2D imagesVal 做分割推理，并用 labelsVal 计算指标。 & boundary & 0.985 & Ext \\
15 & 将私有 2D 数据整理成 DentalClaw canonical case 格式，再导出分割训练格式。 & boundary & 0.985 & Ext \\
16 & 用私有 2D 数据训练牙齿检测模型。 & boundary & 0.985 & Ext \\
17 & 把私有 2D segmentation mask 自动转成检测框并训练 detector。 & boundary & 0.985 & Ext \\
18 & 用私有 2D 牙片训练疾病分类模型。 & boundary & 0.940 & Rej \\
\hline
\multicolumn{5}{@{}l@{}}{\textbf{Ambiguous requests}} \\
19 & 帮我用 TDD 做一个区域分割模型，训练好后告诉我效果。 & ambiguous & 0.868 & Ext \\
20 & 帮我从 TDD 里挑一个病例出一份牙齿相关报告。 & ambiguous & 1.000 & Clar \\
21 & 对 ToothFairy3 images-only CBCT 子集运行已有模型推理，不做监督训练。 & ambiguous & 1.000 & Clar \\
22 & 帮我对 ToothFairy3 做质量分类，看看哪些病例风险比较高。 & ambiguous & 0.868 & Ext \\
23 & 帮我按 ToothFairy3 的子集类型做一个分类实验。 & ambiguous & 0.868 & Ext \\
24 & 帮我给私有 2D 数据出一份临床结果报告。 & ambiguous & 1.000 & Clar \\
25 & 帮我判断私有 2D 里哪些病例的标注能用，哪些需要复核。 & ambiguous & 1.000 & Clar \\
\hline
\multicolumn{5}{@{}l@{}}{\textbf{Trap (out-of-scope / QC-first) requests}} \\
26 & 检查 TDD 检测框是否存在越界、重复或病例泄漏；如果发现阻断性缺陷就停止并列出证据。 & trap & 0.810 & Ext \\
27 & 用 TDD 的图像质量标签训练可用/需复核分类器；训练前先核验质量标签清单，标签不完整就停止。 & trap & 0.810 & Ext \\
28 & 训练 ToothFairy3 前先严格检查标签值是否合法；若有非法标签值就阻断训练并列出病例。 & trap & 1.000 & Clar \\
29 & 先检查私有 2D 数据的 train/val/test split 是否存在泄漏或错配；发现阻断问题则不要训练。 & trap & 1.000 & Clar \\
30 & 将私有 2D 数据导出成 COCO 检测格式；如果缺少检测框标注就停止并报告缺陷。 & trap & 0.810 & Ext \\
\hline
\multicolumn{5}{@{}c@{}}{Overall average: \textbf{0.958} \quad (all intents $\ge$ 0.70 pass threshold)} \\
\end{tabular}

\textbf{Expected-outcome specification.} Expected outcomes were prespecified on the basis of the offline registry contents and the platform's documented scope. Vague requests (entries 19--25) and quality-control-first requests whose text does not state a task family (entries 28 and 29) were prespecified to be answered with conservative clarification, because the platform's parser requests clarification when the task is not stated; execution-level quality-control blocking lies outside this decision-level protocol. Two classification requests (entries 10 and 18) were prespecified to be rejected under the registry scope policy (\texttt{reject\_or\_explain}). These expectations, together with the mapping of every request to its expected outcome, are recorded in \texttt{benchmark\_intents/build\_platform\_mvp\_30.py}.
\end{table}

\subsection{Comparison with general-purpose assistants}

The same 30 requests were issued to three general-purpose assistants: Copilot (GPT-5.6 Luna), Qwen (Qwen3.8-Max), and Claude Code operating with a DeepSeek V4 Pro backend (the same backend model used by the platform's reasoning layer). Each request was issued in a fresh single-turn conversation with an identical preamble describing the available datasets, the six registry entries, the decision rules, and the required JSON output format; the assistants returned a decision without executing any code, and their responses were scored with the same rubric against the same prespecified expected outcomes (Table~S6). The comparison was restricted to the decision layer, each request was issued once, and the assistants' system prompts were not controlled.

\begin{table}[!ht]
\centering
\small
\caption{Decision-evaluation scores of three general-purpose assistants compared with DentalClaw (same 30 requests, same six-dimension rubric).}
\label{tab:supp_assistants}
\begin{tabular}{l c c c c c c}
\toprule
System & Overall & Standard & Boundary & Ambiguous & Trap & Passed \\
\midrule
DentalClaw (platform) & \textbf{0.958} & 1.000 & 0.978 & \textbf{0.943} & 0.886 & 30/30 \\
Copilot (GPT-5.6 Luna) & 0.893 & 1.000 & 0.958 & 0.739 & 0.833 & 28/30 \\
Qwen (Qwen3.8-Max) & 0.891 & 1.000 & 0.958 & 0.739 & 0.818 & 28/30 \\
Claude Code (DeepSeek V4 Pro backend) & 0.889 & 1.000 & 0.954 & 0.739 & 0.818 & 28/30 \\
\bottomrule
\end{tabular}
\end{table}

\subsection{Scoring rules}

Each intent received a score computed as a weighted sum of six prespecified dimensions: planning (0.35), QC blocking (0.20), intent parsing (0.15), external proposal quality (0.15), ambiguity handling (0.10), and boundary identification (0.05). Planning checks whether the \texttt{supported}, \texttt{executable}, and selected-method fields match the prespecified expectation; QC blocking checks whether the request was blocked or released as prespecified; intent parsing checks whether dataset, task family, and mode were parsed correctly; external proposal quality rewards informative non-executable proposals when no registry method matches; ambiguity handling requires ambiguous requests to be blocked; and boundary identification credits correct blocking or half-credits a non-executable external proposal for out-of-scope requests. The overall score is the arithmetic mean of the per-intent scores, with scores at or above 0.70 considered passing.

\section{Manual comparison protocol}

The manual comparison followed the same predefined workflow as the platform runs, without additional hyperparameter tuning or model selection. It was performed by the AI researcher on the study team (the practitioner-level operator of the configuration comparison in Table~S3), who was familiar with the datasets and tooling. Both arms were timed as elapsed wall-clock time on the same NVIDIA RTX 3090 workstation. The comparison is reported as a workflow demonstration rather than a controlled human-factor experiment.

\end{document}
